% from ega4-1-fr.tex
\oldpage[IV-1]{223}

Les sujets traités dans ce Chapitre appellent les remarques suivantes~:

\medskip
\noindent
\emph{a)} Leur caractère commun est d'être relatifs à des propriétés
\emph{locales} de préschémas ou de morphismes, i.e. considérées en un point,
ou aux points d'une fibre, ou dans un voisinage (non spécifié) d'un point ou
d'une fibre. Ces propriétés sont généralement de nature \emph{topologique},
\emph{différentielle} ou \emph{dimensionnelle} (i.e. faisant intervenir les
notions de \emph{dimension} et de \emph{profondeur}) et sont liées aux
propriétés des \emph{anneaux locaux} aux points considérés. Un problème type
est de mettre en relation, pour un morphisme donné $f:X\to Y$ et un point
$x\in X$, les propriétés de $X$ en $x$ avec celles de $Y$ en $y=f(x)$ et de
la fibre $X_y=f^{-1}(y)$ en $x$. Un autre est de déterminer la nature
topologique (par exemple la constructibilité, ou le fait d'être ouvert ou
fermé) de l'ensemble des points $x\in X$ en lesquels $X$ possède une certaine
propriété ou pour lesquels la fibre $X_{f(x)}$ passant par $x$ possède une
certaine propriété en $x$. De même, on s'intéressera à la nature topologique
de l'ensemble des points $y\in Y$ tels que $X$ ait une certaine propriété en
tous les points de la fibre $X_y$, ou tels que cette fibre ait une certaine
propriété.

\medskip
\noindent
\emph{b)} Les notions les plus importantes pour les chapitres suivants sont
celle de \emph{morphisme plat de présentation finie} et celles de
\emph{morphisme lisse} et de \emph{morphisme étale}, qui en sont des cas
particuliers. Leur étude détaillée (jointe à celle de questions connexes)
commence à proprement parler au §~11.

\medskip
\noindent
\emph{c)} Les §§~1 à 10 peuvent être considérés comme de nature préliminaire
et développent trois types de techniques, utilisées non seulement dans les
autres paragraphes du chapitre mais aussi bien entendu dans les chapitres
suivants~:

\medskip
\noindent
\emph{c 1)} Aux §§~1 à 4 sont envisagés divers aspects de la notion de
\emph{changement de base}, surtout en relation avec des conditions de
\emph{finitude} ou de \emph{platitude}~; on y amorce la technique de
\emph{descente} sous ses aspects les plus élémentaires (les questions
d'«~effectivité~» liées à cette technique seront étudiées dans le chap.~V).

\medskip
\noindent
\emph{c 2)} Les §§~5 à 7 sont centrés sur ce que l'on peut appeler des
\emph{techniques noethériennes}, les préschémas considérés étant toujours
supposés localement noethériens, tandis qu'au contraire aucune condition de
finitude n'est généralement imposée aux \emph{morphismes}~; cela tient
essentiellement à ce que les notions de dimension et de profondeur ne sont
guère maniables que pour les anneaux locaux noethériens. Rappelons que le §~7
constitue une théorie «~fine~» des anneaux locaux noethériens, assez peu
utilisée dans la suite du chapitre.

\medskip
\noindent
\emph{c 3)} Les §§~8 à 10 donnent entre autres le moyen d'\emph{éliminer les
hypothèses noethériennes} sur les préschémas étudiés, en leur substituant des
hypothèses convenables de \emph{finitude} (la «~présentation finie~») sur les
\emph{morphismes} considérés~: l'avantage de cette substitution est que ces
dernières hypothèses sont \emph{stables par changement de base}, ce qui n'est
pas le cas pour les hypothèses noethériennes sur les préschémas. La technique
permettant ce passage repose, d'une part, sur l'utilisation de la notion de
\emph{limite projective} de préschémas, grâce à laquelle on peut ramener une
question à la même question sous des hypothèses \emph{noethériennes}~; d'autre
part, sur l'emploi systématique des \emph{ensembles constructibles}, qui ont
le double intérêt de se conserver par image réciproque (par un morphisme

